% =====================================================================
% Guía del estudiante -- Física 9° -- Trimestre I
% Hilo conductor: De Aristóteles a Galileo
% Temas (propuesta de curriculo/plan-anual.md, clases 002-008):
%   cientificos, describir, mru, marcos
% Plan: guia-periodo-I-movimiento.plan.md
% =====================================================================
% Plan del trimestre aprobado por la docente el 2026-09-18. Esta guía es todavía el borrador de
% la generación rápida del 2026-09-14: falta rehacerla con el proceso de tres etapas.
\documentclass[11pt]{article}
\makeatletter\def\input@path{{../../../../plantillas/estilo/}}\makeatother
\usepackage{mmcantillo}

% ======================== DATOS DE LA GUÍA ============================
\tipodocumento{Guía del estudiante}
\asignatura{Física}
\titulo{Describir el movimiento: posición, velocidad y marcos de referencia}
\grado{9°}
\periodo{I}
\hiloconductor{De Aristóteles a Galileo}
% ======================================================================
% DBA: naturales grado 9 · DBA 1 — Comprende que el movimiento de un cuerpo, en un marco de
% referencia inercial dado, se puede describir con gráficos y predecir por medio de
% expresiones matemáticas. (dba/naturales/grados/grado09.tex; en el trimestre I, solo en una
% dimensión y con velocidad constante.)

\begin{document}

\encabezado

% ---------------------------------------------------------------------
% 1. Frase introductoria
% ---------------------------------------------------------------------
% verificado: https://galileoandeinstein.phys.virginia.edu/tns_draft/tns_153to160.html — Crew y de Salvio (1914), pp. 153–154: «My purpose is to set forth a very new science dealing with a very ancient subject. There is, in nature, perhaps nothing older than motion»
\frase{Me propongo exponer una ciencia muy nueva sobre un tema muy antiguo. Quizá no haya
en la naturaleza nada más antiguo que el movimiento [\ldots]}
{Galileo Galilei, \emph{Dos nuevas ciencias} (1638), Jornada tercera; trad. propia de la
versión inglesa de Crew y de Salvio}

% ---------------------------------------------------------------------
% 2. Introducción
% ---------------------------------------------------------------------
\seccion{Introducción}

Todo se mueve: el bus del MIO que te lleva al colegio, el balón en el recreo, la Tierra
alrededor del Sol. Pero ¿qué significa \emph{describir} un movimiento? Decir «va rápido» no
basta. En este trimestre aprenderemos a describirlo como lo hace la física: con una posición
medida desde un punto elegido, con una velocidad que tiene signo, con gráficas y con
ecuaciones que permiten \emph{predecir} dónde estará un cuerpo dentro de un rato.

Nuestro hilo conductor continúa las exposiciones de científicos que ustedes preparan. Durante
casi dos mil años se aceptó la explicación de Aristóteles; en la Edad Media, Nicole Oresme
empezó a dibujar el movimiento; y en el siglo \textsc{xvii} Galileo Galilei lo midió y lo
escribió en lenguaje matemático. En cada tema daremos un paso de ese camino.

\textbf{Al terminar este trimestre podrás:}
\begin{itemize}
  \item distinguir distancia recorrida de desplazamiento, y rapidez de velocidad;
  \item construir e interpretar gráficas posición--tiempo y velocidad--tiempo de un
        movimiento con velocidad constante;
  \item predecir la posición de un cuerpo con la ecuación $x = x_0 + v\,t$;
  \item explicar cómo cambian la velocidad y la gráfica de un movimiento cuando se describe
        desde otro marco de referencia.
\end{itemize}

Este trimestre trabaja el Derecho Básico de Aprendizaje de física de 9.º del Ministerio de
Educación Nacional~\cite{mendba}. En el trimestre I lo estudiamos en una dimensión y con
velocidad constante; la aceleración y el movimiento en el plano vienen en el trimestre II.

\begin{dba}[Grado 9 · DBA 1]
Comprende que el movimiento de un cuerpo, en un marco de referencia inercial dado, se puede
describir con gráficos y predecir por medio de expresiones matemáticas.
\end{dba}

% ---------------------------------------------------------------------
% 3. Marco teórico
% ---------------------------------------------------------------------
\seccion{Marco teórico}

\subsubsection*{Aristóteles: todo movimiento necesita un motor}

En su \emph{Física}, escrita en el siglo \textsc{iv} a.\,C., Aristóteles distinguió dos clases
de movimiento. El \emph{natural}, que un cuerpo realiza por sí mismo según su naturaleza: lo
pesado cae hacia el centro, lo liviano sube. Y el \emph{forzado} o «violento», como el de una
piedra lanzada, que va contra esa naturaleza. Para él, todo movimiento necesita algo que lo
produzca~\cite{sepAristoteles}. Era una explicación del \emph{porqué}; no decía \emph{cuánto}
se mueve un cuerpo ni \emph{cuándo} llega.

\subsubsection*{Oresme: dibujar el movimiento}

En la década de 1350, el francés Nicole Oresme (c.\,1323--1382) escribió un tratado sobre las
«configuraciones» de las cualidades y los movimientos. Allí representó una magnitud que
cambia ---por ejemplo, la velocidad de un cuerpo--- con una figura en la que una línea
representa el tiempo, y con esas figuras demostró geométricamente una regla sobre la
velocidad media~\cite{oresme}. Es un antepasado directo de las gráficas que usaremos en este
trimestre.

\subsubsection*{Galileo: medir}

Galileo Galilei (1564--1642) cambió la pregunta: en lugar de preguntar por qué se mueven los
cuerpos, se preguntó \emph{cómo} se mueven, y lo midió. En sus \emph{Discursos} de 1638
definió el movimiento uniforme como aquel en que un cuerpo recorre distancias iguales en
tiempos iguales, \emph{cualesquiera} que sean esos tiempos~\cite{crew1914}. Y para estudiar la
caída hizo rodar una esfera por un canal inclinado. Como no tenía cronómetro, medía el tiempo
con un reloj de agua: dejaba salir un chorrito de un recipiente durante cada bajada y pesaba
el agua recogida. Encontró que las distancias recorridas en tiempos iguales crecen como los
números impares $1, 3, 5, 7, \ldots$~\cite{galileo1638}

Seguramente has oído que Galileo dejó caer dos bolas desde la torre inclinada de Pisa. Esa
historia viene de una biografía que escribió su secretario, Vincenzo Viviani, en 1654 (se
publicó en 1717); Galileo no dejó ningún registro de ese experimento, y la mayoría de los
historiadores lo consideran una \emph{leyenda}~\cite{physicsworld}. En ciencia importa
distinguir lo documentado de lo que solo se cuenta.

\subsubsection*{El barco de Galileo}

En 1632, en su \emph{Diálogo sobre los dos máximos sistemas del mundo}, Galileo propuso un
experimento imaginario que es la base del último tema de este trimestre~\cite{galileo1632}:

% verificado: https://www.robertoocca.net/fis/fren/fr_principi/galileo_gran_navilio.htm — Dialogo (1632), Jornada segunda: «Riserratevi con qualche amico nella maggiore stanza che sia sotto coverta di alcun gran navilio, e quivi fate d'aver mosche, farfalle e simili animaletti volanti [...] voi non riconoscerete una minima mutazione in tutti li nominati effetti, né da alcuno di quelli potrete comprender se la nave cammina o pure sta ferma»
\begin{cita}{Galileo Galilei, \emph{Diálogo sobre los dos máximos sistemas del mundo} (1632), Jornada segunda, trad. propia~\cite{galileo1632}}
Enciérrense con algún amigo en la cámara más grande que haya bajo cubierta de un gran navío,
y lleven allí moscas, mariposas y otros animalitos voladores [\ldots] no reconocerán el menor
cambio en ninguno de los efectos mencionados, ni podrán saber por ninguno de ellos si la nave
avanza o está quieta.
\end{cita}

Si el barco se mueve con velocidad constante, todo ocurre adentro igual que si estuviera
quieto. El movimiento siempre se describe \emph{respecto a algo}: a eso lo llamamos marco de
referencia.

\subsubsection*{Katherine Johnson: trayectorias calculadas a mano}

Tres siglos después, describir movimientos con precisión llevó a los seres humanos al
espacio. Katherine Johnson (1918--2020) entró en 1953 al centro Langley de la agencia que se
convertiría en la NASA. Hizo el análisis de la trayectoria del vuelo de Alan Shepard, el
primer estadounidense en el espacio (1961), y antes del vuelo orbital de John Glenn (1962)
verificó a mano, con su calculadora de escritorio, las cuentas de la computadora: Glenn pidió
que ella las revisara antes de despegar. En 2015 recibió la Medalla Presidencial de la
Libertad~\cite{nasaJohnson}. Calcular una trayectoria exige, antes que nada, elegir desde
dónde se mide.

% ---------------------------------------------------------------------
% 4. Temas
% ---------------------------------------------------------------------
\tema{Científicos que cambiaron la física}\label{tema:cientificos}

\begin{hilo}
En las exposiciones conocimos personas que cambiaron la manera de entender el mundo. Hay un
cambio que nos interesa especialmente: el paso de \emph{explicar} el movimiento con palabras,
como Aristóteles, a \emph{describirlo} con mediciones, como Galileo.
\end{hilo}

Describir un movimiento es decir \emph{dónde} está un cuerpo en cada instante. Para eso
Galileo necesitó dos cosas que hoy parecen obvias: una manera de medir el tiempo y una
definición precisa. Su definición de movimiento uniforme sigue vigente:

\begin{definicion}[Movimiento uniforme (Galileo, 1638)]
Un movimiento es uniforme cuando el cuerpo recorre distancias iguales en intervalos de tiempo
iguales, cualesquiera que sean esos intervalos.
\end{definicion}

La palabra «cualesquiera» es importante: no basta con que un carro recorra \qty{60}{km} en
cada hora; también debe recorrer \qty{1}{km} en cada minuto y $\frac{1}{60}$~km en cada segundo.
Con su plano inclinado Galileo encontró, en cambio, un movimiento que \emph{no} es uniforme:
en tiempos iguales la esfera recorre distancias que crecen como $1, 3, 5, 7, \ldots$ Ese
movimiento acelerado lo estudiaremos en el trimestre II.

\begin{aplicacion}[Medir distancias con el tiempo]
Durante una tormenta, la luz del relámpago llega casi al instante, pero el sonido del trueno
viaja a unos \qty{340}{m/s}. Contando los segundos entre el relámpago y el trueno se estima a
qué distancia cayó el rayo: es un movimiento uniforme (el del sonido) usado como instrumento de
medida, igual que Galileo usaba el agua de su reloj.
\end{aplicacion}

\begin{preguntas}
% PENDIENTE: el banco no tiene ejercicios verificados de historia del movimiento; solo se usa uno.
% Ejercicio: cinematica-10-033
\pregunta Según una regla práctica, cada cinco segundos entre un relámpago y su trueno indican
una milla de distancia. Suponiendo que la luz llega instantáneamente, estima con esta regla la
rapidez del sonido en m/s. ¿Cuál sería la regla en kilómetros?
\espacio{3cm}
\end{preguntas}

\begin{resumen}
\begin{itemize}
  \item Aristóteles explicaba por qué se mueven los cuerpos; Galileo midió cómo se mueven.
  \item Movimiento uniforme: distancias iguales en tiempos iguales, cualesquiera que sean.
  \item La historia de la torre de Pisa es una leyenda; el plano inclinado y el reloj de agua
        están documentados en los \emph{Discursos} (1638).
\end{itemize}
\end{resumen}

% ---------------------------------------------------------------------
\tema{Describir el movimiento}\label{tema:describir}

\begin{hilo}
Oresme dibujó el tiempo como una línea. Antes de dibujar, hay que decidir desde dónde medimos
y hacia dónde contamos positivo: hay que elegir un marco de referencia.
\end{hilo}

\begin{definicion}[Marco de referencia y posición]
Un \textbf{marco de referencia} es un cuerpo u objeto desde el que se mide, con un origen y un
sentido positivo. En una dimensión, la \textbf{posición} $x$ de un cuerpo es su coordenada en
ese eje: positiva a un lado del origen y negativa al otro.
\end{definicion}

La \textbf{trayectoria} es el camino que sigue el cuerpo. La \textbf{distancia recorrida} es la
longitud total de ese camino (siempre positiva). El \textbf{desplazamiento} es el cambio de
posición, $\Delta x = x_{\text{final}} - x_{\text{inicial}}$: tiene signo y solo depende de dónde
empezó y dónde terminó el movimiento~\cite{recurso10}.

\begin{definicion}[Rapidez media y velocidad media]
\[
  \text{rapidez media} = \frac{\text{distancia recorrida}}{\text{tiempo}},
  \qquad
  \bar v = \frac{\Delta x}{\Delta t} = \frac{x_2 - x_1}{t_2 - t_1}.
\]
La rapidez media es siempre positiva; el signo de la velocidad media indica el sentido del
movimiento.
\end{definicion}

\begin{ejemplo}[Una caminata de ida y vuelta]
Una persona camina \qty{70}{m} hacia el oriente y luego \qty{30}{m} hacia el occidente, en
\qty{70}{s} en total. La distancia recorrida es \qty{100}{m}, pero el desplazamiento es
$\Delta x = \qty{+40}{m}$. Entonces la rapidez media es
$100/70 \approx \qty{1,4}{m/s}$ y la velocidad media, $40/70 \approx \qty{0,57}{m/s}$ hacia el
oriente~\cite{recurso10}.
\end{ejemplo}

\begin{ejemplo}[Velocidad negativa]
Un corredor pasa de $x_1 = \qty{50,0}{m}$ a $x_2 = \qty{30,5}{m}$ en \qty{3,00}{s}:
$\bar v = (30{,}5 - 50{,}0)/3{,}00 = \qty{-6,50}{m/s}$. El signo menos indica que se mueve
hacia el lado negativo del eje~\cite{recurso10}.
\end{ejemplo}

Para pasar de km/h a m/s se multiplica por $\frac{1000\ \text{m}}{1\ \text{km}}$ y por
$\frac{1\ \text{h}}{3600\ \text{s}}$, es decir, se divide entre $3{,}6$:
$\qty{36}{km/h} = \qty{10}{m/s}$.

\begin{aplicacion}[El odómetro y el GPS]
El odómetro de un carro o de una bicicleta mide la distancia recorrida: sigue sumando aunque
des vueltas por el mismo barrio. La aplicación de mapas de tu celular, en cambio, puede
decirte a qué distancia en línea recta estás del punto de partida, que es la magnitud de tu
desplazamiento. Al final de un paseo que vuelve a la casa, el odómetro marca kilómetros y el
desplazamiento es cero.
\end{aplicacion}

\begin{preguntas}
% Ejercicio: cinematica-10-001
\pregunta Una hormiga parte de $x = 20$ cm sobre una hoja cuadriculada y camina a lo largo del
eje $x$ hasta $x = -20$ cm. Luego se devuelve y camina hasta $x = -10$ cm. ¿Cuál es el
desplazamiento de la hormiga y cuál la distancia total recorrida?
\espacio{2.5cm}
% Ejercicio: cinematica-10-030
\pregunta ¿Cuál debe ser la rapidez promedio de tu automóvil para recorrer 235 km en
$\num{3,25}$ h?
\espacio{2cm}
% Ejercicio: cinematica-10-036
\pregunta Un caballo se aleja de su entrenador galopando en línea recta 116 m en $\num{14,0}$ s.
Luego se devuelve y recorre la mitad de esa distancia en $\num{4,8}$ s. Calcula su rapidez
promedio.
\espacio{2.5cm}
% Ejercicio: cinematica-10-037
\pregunta Para el caballo anterior (116 m alejándose en $\num{14,0}$ s y 58 m de regreso en
$\num{4,8}$ s), calcula la velocidad promedio de todo el viaje, tomando «alejándose del
entrenador» como sentido positivo.
\espacio{2.5cm}
% Ejercicio: cinematica-10-022
\pregunta Viajas del punto A al B con rapidez constante de 70 km/h y luego la misma distancia de
B a C con rapidez constante de 90 km/h. ¿Tu rapidez promedio para el viaje completo es de
80 km/h? Explica.
\renglones{4}
\end{preguntas}

\begin{resumen}
\begin{itemize}
  \item Toda posición se mide desde un marco de referencia con origen y sentido positivo.
  \item Distancia: longitud del camino. Desplazamiento: $\Delta x = x_2 - x_1$, con signo.
  \item Rapidez media = distancia/tiempo; velocidad media = desplazamiento/tiempo.
  \item De km/h a m/s se divide entre $3{,}6$.
\end{itemize}
\end{resumen}

% ---------------------------------------------------------------------
\tema{Movimiento rectilíneo uniforme}\label{tema:mru}

\begin{hilo}
Galileo definió el movimiento uniforme con palabras; Oresme habría sabido dibujarlo. Juntemos
las dos ideas: una ecuación y una gráfica para el mismo movimiento.
\end{hilo}

\begin{definicion}[Movimiento rectilíneo uniforme (MRU)]
Un cuerpo tiene MRU cuando se mueve en línea recta con velocidad constante $v$. Si en $t = 0$
está en $x_0$, su posición en cualquier instante es
\[
  x = x_0 + v\,t.
\]
\end{definicion}

En la gráfica posición--tiempo ($x$--$t$) un MRU es una \textbf{recta}: su punto de corte con
el eje vertical es $x_0$ y su \textbf{pendiente} es la velocidad. Una pendiente negativa
significa que el cuerpo se mueve hacia el lado negativo del eje. En la gráfica
velocidad--tiempo ($v$--$t$) un MRU es una recta horizontal, y el \textbf{área} bajo ella es el
desplazamiento.

\begin{ejemplo}[Dos movimientos en la misma gráfica]
Un carro de juguete sigue $x = 20 + 5t$ (en metros y segundos): en $t = 0, 2, 4, 6$ s está en
$20, 30, 40, 50$ m. Un perro que regresa a su casa sigue $x = 60 - 10t$: llega al origen en
$t = \qty{6}{s}$. En la gráfica de la izquierda, la recta del carro sube (velocidad
$+\qty{5}{m/s}$) y la del perro baja (velocidad $\qty{-10}{m/s}$). En la de la derecha, el
área sombreada es el desplazamiento del carro en 6 s: $5 \times 6 = \qty{30}{m}$, que es
justamente $50 - 20$.

\begin{center}
\begin{tikzpicture}
\begin{axis}[mmc, width=0.47\linewidth, height=0.4\linewidth,
             xmin=0, xmax=7, ymin=0, ymax=65, xlabel={$t$ (s)}, ylabel={$x$ (m)},
             xtick={0,2,4,6}, ytick={0,20,40,60},
             legend pos=north east, legend style={draw=none, font=\footnotesize}]
  \addplot+[domain=0:6.5, samples=2, mark=none] {20+5*x};
  \addplot+[domain=0:6, samples=2, mark=none] {60-10*x};
  \legend{carro, perro}
\end{axis}
\end{tikzpicture}\hfill
\begin{tikzpicture}
\begin{axis}[mmc, width=0.47\linewidth, height=0.4\linewidth,
             xmin=0, xmax=7, ymin=0, ymax=7, xlabel={$t$ (s)}, ylabel={$v$ (m/s)},
             xtick={0,2,4,6}, ytick={0,5}]
  \fill[black!15] (axis cs:0,0) rectangle (axis cs:6,5);
  \addplot+[domain=0:6.5, samples=2, mark=none] {5};
  \node[font=\footnotesize] at (axis cs:3,2.5) {área $= \qty{30}{m}$};
\end{axis}
\end{tikzpicture}
\end{center}
\end{ejemplo}

\begin{aplicacion}[Distancias de seguridad]
Mirar el celular dos segundos mientras se conduce parece poco. Pero a velocidad constante el
carro sigue avanzando $v\,t$ metros sin que nadie lo controle: a \qty{110}{km/h} son más de
sesenta metros. Por eso las normas de tránsito piden dejar una distancia de seguridad que
crece con la velocidad.
\end{aplicacion}

\begin{preguntas}
% PENDIENTE (evidencia 1): no hay en el banco un ejercicio verificado de leer una gráfica x–t y
% escribir x = x0 + v t (cinematica-10-038/039 citan una imagen que la guía no tiene).
% Ejercicio: cinematica-10-029
\pregunta Si vas manejando a 110 km/h por una carretera recta y te distraes durante
$\num{2,0}$ s, ¿qué distancia recorres sin prestar atención?
\espacio{2.5cm}
% Ejercicio: cinematica-10-043
\pregunta Una bola de boliche que rueda con rapidez constante golpea los pinos al final de una
pista de $\num{16,5}$ m. El jugador oye el golpe $\num{2,50}$ s después de lanzarla. ¿Cuál es la
rapidez de la bola, si la del sonido es 340 m/s?
\espacio{3cm}
% Ejercicio: cinematica-10-040
\pregunta En un CD de audio, los bits se codifican a lo largo de una espiral y cada bit ocupa
aproximadamente $\num{0,28}\ \mu$m. El lector láser recorre la espiral a una rapidez constante
de aproximadamente $\num{1,2}$ m/s. Determina el número $N$ de bits que el reproductor lee cada
segundo.
\espacio{2.5cm}
\end{preguntas}

\begin{resumen}
\begin{itemize}
  \item MRU: $x = x_0 + v\,t$, con $v$ constante.
  \item Gráfica $x$--$t$: una recta; su pendiente es la velocidad y su corte, $x_0$.
  \item Gráfica $v$--$t$: una recta horizontal; el área bajo ella es el desplazamiento.
  \item Con la ecuación se predice dónde estará un cuerpo y cuándo se encuentran dos cuerpos.
\end{itemize}
\end{resumen}

% ---------------------------------------------------------------------
\tema{Marcos de referencia}\label{tema:marcos}

\begin{hilo}
Volvamos al barco de Galileo. Para quien está en el camarote, una gota que cae del techo cae
en línea recta; para alguien en la orilla, la gota además avanza con el barco. ¿Quién tiene
la razón? Los dos: cada uno describe el movimiento desde su propio marco.
\end{hilo}

Si un cuerpo $A$ tiene velocidad $v_A$ y un observador $B$ se mueve con velocidad $v_B$ (las
dos medidas desde el suelo y en la misma recta), el observador $B$ ve moverse a $A$ con la
\textbf{velocidad relativa}
\[
  v_{A\text{ respecto a }B} = v_A - v_B.
\]
Al cambiar de marco cambian la velocidad, la ecuación de posición y la gráfica; lo que no
cambia es \emph{lo que ocurre}: si dos cuerpos se chocan en un marco, se chocan en todos.

\begin{ejemplo}[Un pasajero en el bus]
Un bus va a \qty{12}{m/s} y un pasajero camina hacia el frente a \qty{1,5}{m/s} respecto al
bus. Respecto a la calle, el pasajero va a $12 + 1{,}5 = \qty{13,5}{m/s}$; si camina hacia
atrás, a $12 - 1{,}5 = \qty{10,5}{m/s}$ (todavía hacia adelante). En la gráfica, la misma
caminata hacia el frente vista desde la calle ($x = 13{,}5\,t$) y desde el bus ($x' = 1{,}5\,t$):
dos rectas con distinta pendiente para un solo movimiento.

\begin{center}
\begin{tikzpicture}
\begin{axis}[mmc, width=0.6\linewidth, height=0.4\linewidth,
             xmin=0, xmax=4.5, ymin=0, ymax=60, xlabel={$t$ (s)}, ylabel={posición (m)},
             xtick={0,1,2,3,4}, ytick={0,20,40},
             legend pos=outer north east, legend style={draw=none, font=\footnotesize}]
  \addplot+[domain=0:4.2, samples=2, mark=none] {13.5*x};
  \addplot+[domain=0:4.2, samples=2, mark=none] {1.5*x};
  \legend{desde la calle, desde el bus}
\end{axis}
\end{tikzpicture}
\end{center}
\end{ejemplo}

El ejemplo del propio DBA va un paso más allá: un pasajero de un bus que va con velocidad
constante deja caer una pelota. Para él, la pelota cae en línea recta vertical; para un peatón
en la acera, la misma pelota describe una curva, porque además avanza con el bus. Esa curva,
la parábola, la estudiaremos en el trimestre II.

\begin{aplicacion}[Servir el café en un avión]
En un avión en vuelo de crucero, a unos \qty{900}{km/h} respecto al suelo, la azafata sirve
el café en el vaso sin apuntar «hacia atrás», y el café cae dentro. Es el barco de Galileo en
versión moderna: mientras la velocidad es constante, dentro del avión todo ocurre como si
estuviera quieto. Solo al despegar, frenar o atravesar turbulencia (cuando la velocidad
cambia) se nota el movimiento.
\end{aplicacion}

\begin{preguntas}
% PENDIENTE (evidencia 3): no hay en el banco un ejercicio verificado de cambiar la gráfica o la
% ecuación de un movimiento a otro marco; solo el de velocidad relativa.
% Ejercicio: cinematica-10-042
\pregunta Un automóvil que viaja a 95 km/h va 110 m detrás de un camión que viaja a 75 km/h.
¿Cuánto tiempo le tomará al automóvil alcanzar al camión?
\espacio{3cm}
% Pendiente de aprobación de la docente (manual-pendiente), no se incluye todavía:
% Ejercicio: cinematica-2d-10-021 (persona en un vagón de tren que lanza una pelota)
\end{preguntas}

\begin{resumen}
\begin{itemize}
  \item La descripción de un movimiento depende del marco de referencia.
  \item Velocidad relativa en una recta: $v_{A\text{ resp. }B} = v_A - v_B$.
  \item Al cambiar de marco cambian la ecuación y la pendiente de la gráfica $x$--$t$.
  \item Con velocidad constante, ningún experimento interior distingue el reposo del
        movimiento (el barco de Galileo).
\end{itemize}
\end{resumen}

% ---------------------------------------------------------------------
% 5. Saber 11
% ---------------------------------------------------------------------
\seccion{Prepárate para Saber 11}

\begin{preguntas}
% PENDIENTE: el banco no tiene preguntas tipo Saber de cinemática con contexto colombiano.
% Ejercicio: cinematica-10-002
\pregunta Un automóvil viaja a una rapidez constante de 50 km/h durante 100 km. Luego acelera a
100 km/h y recorre otros 100 km. ¿Cuál es la rapidez promedio de su viaje de 200 km?
\begin{opciones*}(4)
  \item 67 km/h
  \item 75 km/h
  \item 81 km/h
  \item 50 km/h
\end{opciones*}

% Ejercicio: cinematica-10-008
\pregunta Un automóvil parte del reposo y acelera a $10\ \mathrm{m/s^2}$ constantes durante una
carrera de un cuarto de milla (402 m). ¿Qué tan rápido viaja cuando cruza la línea de meta?
\begin{opciones*}(4)
  \item 8090 m/s
  \item 90 m/s
  \item 81 m/s
  \item 809 m/s
\end{opciones*}

% Ejercicio: mediciones-10-071
\pregunta Tres estudiantes obtienen las siguientes ecuaciones, donde $x$ es la distancia
recorrida, $v$ la rapidez, $a$ la aceleración, $t$ el tiempo y el subíndice 0 indica el valor
en $t = 0$. ¿Cuál es correcta según una comprobación dimensional?
\begin{opciones*}(3)
  \item $x = vt^{2} + 2at$
  \item $x = v_{0}t + \frac{1}{2}at^{2}$
  \item $x = v_{0}t + 2at^{2}$
\end{opciones*}
\end{preguntas}

% ---------------------------------------------------------------------
% 6. Cierre del hilo conductor
% ---------------------------------------------------------------------
\seccion{Cierre: de explicar a medir}

Aristóteles dio una explicación que duró casi dos mil años. Oresme se atrevió a dibujar el
tiempo. Galileo midió con agua y con un plano inclinado, definió el movimiento uniforme y
entendió que el movimiento se describe siempre desde un marco. Con esas herramientas ---una
posición, una velocidad con signo, una gráfica y una ecuación--- ya puedes predecir dónde
estará un bus del MIO dentro de un minuto.

Pero Galileo dejó una pregunta abierta en su propio plano inclinado: ¿qué pasa cuando la
velocidad \emph{cambia}? En el trimestre II estudiaremos la aceleración, la caída libre y los
movimientos en el plano, y volveremos a la pelota del bus para ver por qué, desde la acera,
describe una parábola.

% ---------------------------------------------------------------------
% 7. Evaluación
% ---------------------------------------------------------------------
\matrizevaluacion
  {Distingue distancia de desplazamiento y rapidez de velocidad, y explica qué es un marco de referencia.}
  {Construye e interpreta gráficas $x$--$t$ y $v$--$t$ de un MRU y predice posiciones con $x = x_0 + v\,t$.}
  {Distingue lo documentado de las leyendas y revisa las unidades de sus resultados.}
  {Expone y escucha con respeto en las presentaciones de científicos y trabaja en equipo en los talleres.}

\autoevaluacion

\seguimientodocente

% ---------------------------------------------------------------------
% 8. Referencias
% ---------------------------------------------------------------------
\begin{referencias}
% verificado: https://archive.org/details/dialoguesconcern00galiuoft — Galileo, Dialogues concerning two new sciences, trad. Crew y de Salvio, Nueva York: Macmillan, 1914; pasaje en pp. 153–154 (https://galileoandeinstein.phys.virginia.edu/tns_draft/tns_153to160.html)
\bibitem{crew1914} Galilei, G. (1914). \emph{Dialogues concerning two new sciences}
  (H. Crew y A. de Salvio, trads.). Nueva York: Macmillan. (Obra original de 1638.)
% verificado: https://portalegalileo.museogalileo.it/igjr.asp?c=36307 — Museo Galileo: Dialogo sopra i due massimi sistemi del mondo, Florencia, Landini, 1632
\bibitem{galileo1632} Galilei, G. (1632). \emph{Dialogo sopra i due massimi sistemi del
  mondo, tolemaico e copernicano}. Florencia: G.\,B. Landini.
% verificado: https://portalegalileo.museogalileo.it/igjr.asp?c=36308 — Museo Galileo: Discorsi e dimostrazioni matematiche intorno a due nuove scienze, Leiden, Elzevir, 1638 (plano inclinado y reloj de agua: https://galileo.library.rice.edu/lib/student_work/experiment95/inclined_plane.html)
\bibitem{galileo1638} Galilei, G. (1638). \emph{Discorsi e dimostrazioni matematiche intorno
  a due nuove scienze}. Leiden: Elzevir.
% verificado: recursos/fisica/Guías pedagógicas Física/markdown/10 - Guia_de_apoyo_cinematica_y_dinamica_grado_10_fisica.md (archivo local), cap. 2
\bibitem{recurso10} Guía de apoyo de Física 10.º: \emph{Cinemática y dinámica}, cap. 2
  (Descripción del movimiento: cinemática en una dimensión). Material del colegio.
% verificado: https://www.colombiaaprende.edu.co/sites/default/files/files_public/2022-06/DBA_C.Naturales-min.pdf — PDF oficial del MEN publicado en Colombia Aprende
\bibitem{mendba} Ministerio de Educación Nacional (2016). \emph{Derechos Básicos de
  Aprendizaje V.1: Ciencias Naturales}. Bogotá: MEN.
% verificado: https://www.nasa.gov/centers-and-facilities/langley/katherine-johnson-biography/ — 1918–2020; Langley desde 1953; Shepard 1961; Glenn 1962; Medalla Presidencial de la Libertad 2015
\bibitem{nasaJohnson} NASA. \emph{Katherine Johnson biography}. Consultado el 14 de
  septiembre de 2026.
% verificado: https://mathshistory.st-andrews.ac.uk/Biographies/Oresme/ — 1323–1382; gráficas de una magnitud variable; prueba de la regla de Merton (década de 1350: https://www.encyclopedia.com/people/philosophy-and-religion/other-religious-beliefs-biographies/nicole-oresme)
\bibitem{oresme} MacTutor History of Mathematics. \emph{Nicole Oresme}. University of St
  Andrews. Consultado el 14 de septiembre de 2026.
% verificado: https://physicsworld.com/a/the-legend-of-the-leaning-tower/ — Viviani es la única fuente; la mayoría de los historiadores dudan (Viviani 1654, publicado 1717: https://upittpress.org/wp-content/uploads/2018/07/9780822944072exr.pdf)
\bibitem{physicsworld} Physics World. \emph{The legend of the leaning tower}. Consultado el 14
  de septiembre de 2026.
% verificado: https://plato.stanford.edu/entries/aristotle-natphil/ — §5: movimiento natural y forzado; todo movimiento necesita un motor
\bibitem{sepAristoteles} Stanford Encyclopedia of Philosophy. \emph{Aristotle's natural
  philosophy}. Consultado el 14 de septiembre de 2026.
\end{referencias}

\end{document}
